% ===========================================================================
%  sviluppatore/08-qualita/03-rilascio.tex
% ===========================================================================
\section{Impacchettare e rilasciare}\label{sez:rilascio}

\subsection{Un comando solo}

\begin{codiceInterno}{console}{Dalla radice del progetto}{build/build\_app.py}
python build/build_app.py                # backend + app per il sistema corrente
python build/build_app.py --backend-only # solo l'eseguibile del backend
python build/build_app.py --app-only     # solo Electron (backend gia' pronto)
\end{codiceInterno}

\noindent
Tre passi, in ordine: PyInstaller produce un backend autonomo che non
richiede Python sulla macchina di chi riceve l'app; il risultato viene
copiato dove electron-builder lo cerca; electron-builder produce
l'installatore per il sistema corrente (\file{.exe}, \file{.dmg},
\file{AppImage}).

\subsection{Dove finiscono i dati dell'utente}

Non dentro l'applicazione. Il bundle è di sola lettura e ogni aggiornamento
lo sovrascrive: al primo avvio il backend crea uno spazio di lavoro nella
cartella utente e vi copia le librerie iniziali --- strategie, approcci,
modalità di pesatura. È lo stesso meccanismo di \simbolo{PROJECT\_ROOT}
descritto nel \vediSez{dati}, visto dall'altro capo.

\begin{attenzione}{Le librerie iniziali si copiano una volta}
Al primo avvio, non a ogni aggiornamento. Un artefatto nuovo aggiunto a una
versione nuova dell'applicazione \emph{non} compare automaticamente nello
spazio di lavoro di chi ha già installato: quello spazio è suo, e
sovrascriverlo cancellerebbe le sue modifiche. È una conseguenza voluta del
punto precedente, e va tenuta presente quando si aggiunge un artefatto di
serie.
\end{attenzione}

\subsection{Gli aggiornamenti automatici}

Chi ha già installato riceve le versioni nuove da solo: all'avvio
\pkg{electron-updater} interroga le release di GitHub, scarica in silenzio e
mostra la notifica «aggiornamento pronto» nella campanella. L'utente riavvia
quando gli fa comodo. Il repository si ricava dal remote git, o si dichiara
con \cod{GH\_REPO}.

\section{Il manifesto delle query}\label{sez:manifesto}

È il meccanismo di sicurezza più importante fra la metà locale e il sito, e
merita di essere capito prima di toccare qualunque cosa interroghi il
database.

\figuraTikz{manifest-query}
  {Il rapporto di fiducia, prima e dopo.}
  {manifest-query}

\begin{codiceInterno}{python}{Il motivo, dal generatore}{build/gen\_query\_manifest.py}
Perche' esiste
--------------
Finora l'app inviava SQL libero a ``/api/dashboard/query``, protetto da un
segreto uguale per tutte le installazioni. Chi estraeva quel segreto poteva
eseguire qualsiasi comando sul database di produzione.

Con questo file il rapporto si rovescia: il sito conosce in anticipo l'elenco
esatto delle query ammesse e l'app puo' solo chiederne una per NOME (un hash
del testo SQL). Qualunque altra cosa viene rifiutata, segreto o no.
\end{codiceInterno}

\subsection{L'identificativo di una query}

\cod{op\_id} è l'impronta SHA-256 del testo SQL normalizzato negli spazi,
troncata a \num{32} cifre esadecimali. Conseguenza diretta: \textbf{cambiare
una query cambia il suo identificativo}. Una query modificata e non
rigenerata nel manifesto viene rifiutata, segreto o no.

\subsection{La legatura all'utente}

Limitare le query non basta: chi è autenticato potrebbe chiedere i dati di un
altro passando un identificativo diverso. Per ogni query il generatore
individua il parametro che rappresenta l'utente, in due forme:

\begin{codiceInterno}{sql}{Le due forme riconosciute}{build/gen\_query\_manifest.py}
    WHERE "userId" = %s              -> conta i %s che lo precedono
    INSERT INTO t(id, "userId", ...) -> posizione nella lista delle colonne
\end{codiceInterno}

\noindent
Il sito \emph{sovrascrive} quel parametro con l'identificativo ricavato dal
token. Anche mentendo, il client ottiene solo i propri dati.

\begin{pratica}{Quando si tocca una query}
\begin{enumerate}
  \item Modificare l'SQL.
  \item \file{python build/gen\_query\_manifest.py} --- riscrive il JSON nel
        sito.
  \item \file{python build/gen\_query\_manifest.py --check} --- verifica che
        sia aggiornato. È il comando da mettere in una verifica automatica.
  \item Rilasciare \emph{prima} il sito, poi l'app: un'app che chiede un
        \cod{op\_id} che il sito non conosce ancora viene rifiutata.
\end{enumerate}
\end{pratica}

\begin{attenzione}{L'ordine di rilascio non è simmetrico}
Il sito accetta sia gli identificativi vecchi sia quelli nuovi finché il
manifesto li contiene entrambi; l'app no. Aggiornare l'app prima del sito
rompe le installazioni appena aggiornate, e il sintomo è un rifiuto senza
spiegazione --- perché dal punto di vista del sito è esattamente ciò che
deve succedere a una query sconosciuta.
\end{attenzione}
